\uuid{SGe5}
\exo7id{3867}
\titre{exo7 3867}
\auteur{quercia}
\organisation{exo7}
\datecreate{2010-03-11}
\isIndication{false}
\isCorrection{false}
\chapitre{Continuité, limite et étude de fonctions réelles}
\sousChapitre{Continuité : théorie}
\cpge{mpsi-mp2i.limites-continuite.b, pcsi.limites-continuite.b, ptsi.limites-continuite.b, mpsi-mp2i.limites-continuite.c, pcsi.limites-continuite.c, ptsi.limites-continuite.c, mpsi-mp2i.fonctions-valeurs-reelles-complexes.a, pcsi.techniques-calcul-fonctions.a, ptsi.techniques-calcul-fonctions.a, bcpst1.limites-continuite.a, bcpst1.limites-continuite.c, bcpst1.fonctions-reelles-usuelles.a}
\module{Analyse}
\niveau{L1}
\difficulte{}

\contenu{
\texte{
Soit $f : {\R} \to {\R}$ telle que : $\forall\ x,y \in \R,\ f(x+y) = f(x) + f(y)$.
On pose $a = f(1)$.
}
\begin{enumerate}
    \item \question{Montrer que : $\forall\ x\in\Q$, $f(x) = ax$.}
    \item \question{On suppose $f$ continue. Montrer que : $\forall\ x \in \R$, $f(x) = ax$.}
    \item \question{On suppose que $f$ est bornée au voisinage de 0.
    Montrer que : $\forall\ x \in \R$, $f(x) = ax$.}
\end{enumerate}
}
